% answers.tex —— 2014 年·全国大纲卷·理科（广西） · 参考答案与详细解析片段
% 由 tools/generate_projects.py 依来源总账自 2014/2014_数学真题_大纲理_试题.md + 2014/2014_数学真题_大纲理_解析正文.md + 2014/2014_数学真题_大纲理_答案要点.md 生成，勿手改（改源 MD 后重跑）


\begin{answer}
\gksolhead{1}{}

\textbf{题目}：

设 $z=\dfrac{10\mathrm i}{3+\mathrm i}$，则 $z$ 的共轭复数为\mbox{（　　）}

（A）$-1+3\mathrm i$；（B）$-1-3\mathrm i$；（C）$1+3\mathrm i$；（D）$1-3\mathrm i$。

\Needspace{4\baselineskip}
\textbf{解析}

\Needspace{5\baselineskip}
分子分母同乘 $3-\mathrm i$，分母化为实数：
\nopagebreak[4]
\[\begin{aligned}
z&=\frac{10\mathrm i}{3+\mathrm i}
=\frac{10\mathrm i(3-\mathrm i)}{(3+\mathrm i)(3-\mathrm i)}\\
&=\frac{10\mathrm i(3-\mathrm i)}{10}
=\mathrm i(3-\mathrm i)\\
&=3\mathrm i-\mathrm i^2=1+3\mathrm i.
\end{aligned}\]
所以 $\bar z=1-3\mathrm i$。

选 D。

\gkanswer{答案：D。$z=\dfrac{10\mathrm i(3-\mathrm i)}{10}=1+3\mathrm i$，共轭复数为 $1-3\mathrm i$。}

\end{answer}
\begin{answer}
\gksolhead{2}{}

\textbf{题目}：

设集合 $M=\{x\mid x^2-3x-4<0\}$，$N=\{x\mid 0\leqslant x\leqslant 5\}$，则 $M\cap N=$\mbox{（　　）}

（A）$(0,4]$；（B）$[0,4)$；（C）$[-1,0)$；（D）$(-1,0]$。

\Needspace{4\baselineskip}
\textbf{解析}

\Needspace{5\baselineskip}
先解二次不等式。$x^2-3x-4=(x-4)(x+1)$，开口向上，故
\nopagebreak[4]
\[x^2-3x-4<0\iff -1<x<4,\]
即 $M=(-1,4)$。又 $N=[0,5]$。

\Needspace{5\baselineskip}
$M$ 与 $N$ 的公共部分为
\nopagebreak[4]
\[M\cap N=[0,4).\]
左端点 $0$ 属于两者，右端点 $4$ 不属于 $M$。

选 B。

\gkanswer{答案：B。$M=(-1,4)$，$N=[0,5]$，故 $M\cap N=[0,4)$。}

\end{answer}
\begin{answer}
\gksolhead{3}{}

\textbf{题目}：

设 $a=\sin33^\circ$，$b=\cos55^\circ$，$c=\tan35^\circ$，则\mbox{（　　）}

（A）$a>b>c$；（B）$b>c>a$；（C）$c>b>a$；（D）$c>a>b$。

\Needspace{4\baselineskip}
\textbf{解析}

\Needspace{5\baselineskip}
先统一成正弦。由余角关系
\nopagebreak[4]
\[b=\cos55^\circ=\sin35^\circ.\]
\Needspace{5\baselineskip}
正弦在 $\left(0,\dfrac\pi2\right)$ 上递增，而 $35^\circ>33^\circ$，故
\nopagebreak[4]
\[b=\sin35^\circ>\sin33^\circ=a.\]
\Needspace{5\baselineskip}
再由
\nopagebreak[4]
\[c=\tan35^\circ=\frac{\sin35^\circ}{\cos35^\circ}>\sin35^\circ=b,\]
\Needspace{5\baselineskip}
因为 $0<\cos35^\circ<1$。于是
\nopagebreak[4]
\[c>b>a.\]

选 C。

\gkanswer{答案：C。$b=\cos55^\circ=\sin35^\circ$，而 $\tan35^\circ>\sin35^\circ>\sin33^\circ$，故 $c>b>a$。}

\end{answer}
\begin{answer}
\gksolhead{4}{}

\textbf{题目}：

若向量 $\boldsymbol a$、$\boldsymbol b$ 满足：$|\boldsymbol a|=1$，$(\boldsymbol a+\boldsymbol b)\perp\boldsymbol a$，$(2\boldsymbol a+\boldsymbol b)\perp\boldsymbol b$，则 $|\boldsymbol b|=$\mbox{（　　）}

（A）$2$；（B）$\sqrt2$；（C）$1$；（D）$\dfrac{\sqrt2}{2}$。

\Needspace{4\baselineskip}
\textbf{解析}

\Needspace{5\baselineskip}
由 $(\boldsymbol a+\boldsymbol b)\perp\boldsymbol a$ 得
\nopagebreak[4]
\[(\boldsymbol a+\boldsymbol b)\cdot\boldsymbol a=|\boldsymbol a|^2+\boldsymbol a\cdot\boldsymbol b=1+\boldsymbol a\cdot\boldsymbol b=0,\]
\Needspace{5\baselineskip}
故
\nopagebreak[4]
\[\boldsymbol a\cdot\boldsymbol b=-1.\]
\Needspace{5\baselineskip}
再由 $(2\boldsymbol a+\boldsymbol b)\perp\boldsymbol b$ 得
\nopagebreak[4]
\[(2\boldsymbol a+\boldsymbol b)\cdot\boldsymbol b=2\boldsymbol a\cdot\boldsymbol b+|\boldsymbol b|^2=-2+|\boldsymbol b|^2=0,\]
故 $|\boldsymbol b|^2=2$，即 $|\boldsymbol b|=\sqrt2$。

选 B。

\gkanswer{答案：B。由 $(\boldsymbol a+\boldsymbol b)\perp\boldsymbol a$ 得 $\boldsymbol a\cdot\boldsymbol b=-1$，再由 $(2\boldsymbol a+\boldsymbol b)\perp\boldsymbol b$ 得 $|\boldsymbol b|^2=2$。}

\end{answer}
\begin{answer}
\gksolhead{5}{}

\textbf{题目}：

有 6 名男医生、5 名女医生，从中选出 2 名男医生、1 名女医生组成一个医疗小组，则不同的选法共有\mbox{（　　）}

（A）60 种；（B）70 种；（C）75 种；（D）150 种。

\Needspace{4\baselineskip}
\textbf{解析}

\Needspace{5\baselineskip}
分两步选取，两个步骤相互独立，用乘法原理：
\nopagebreak[4]
\[\mathrm C_6^2\cdot\mathrm C_5^1=\frac{6\times5}{2}\times5=15\times5=75.\]
即共有 $75$ 种不同选法。

选 C。

\gkanswer{答案：C。$\mathrm C_6^2\mathrm C_5^1=15\times5=75$。}

\end{answer}
\begin{answer}
\gksolhead{6}{}

\textbf{题目}：

已知椭圆 $C:\dfrac{x^2}{a^2}+\dfrac{y^2}{b^2}=1$（$a>b>0$）的左、右焦点分别为 $F_1$、$F_2$，离心率为 $\dfrac{\sqrt3}{3}$，过 $F_2$ 的直线 $l$ 交 $C$ 于 $A$、$B$ 两点，若 $\triangle AF_1B$ 的周长为 $4\sqrt3$，则 $C$ 的方程为\mbox{（　　）}

（A）$\dfrac{x^2}{3}+\dfrac{y^2}{2}=1$；（B）$\dfrac{x^2}{3}+y^2=1$；（C）$\dfrac{x^2}{12}+\dfrac{y^2}{8}=1$；（D）$\dfrac{x^2}{12}+\dfrac{y^2}{4}=1$。

\Needspace{4\baselineskip}
\textbf{解析}

\Needspace{5\baselineskip}
由椭圆定义，$|AF_1|+|AF_2|=2a$，$|BF_1|+|BF_2|=2a$。因为 $A$、$B$ 都在过 $F_2$ 的直线上，且 $F_2$ 在 $A$、$B$ 之间，所以
\nopagebreak[4]
\[\begin{aligned}
\triangle AF_1B\ \text{\text{的周长}}
&=|AB|+|AF_1|+|BF_1|\\
&=(|AF_2|+|BF_2|)+|AF_1|+|BF_1|\\
&=4a.
\end{aligned}\]
\Needspace{5\baselineskip}
由 $4a=4\sqrt3$ 得
\nopagebreak[4]
\[a=\sqrt3.\]
\Needspace{5\baselineskip}
又离心率
\nopagebreak[4]
\[e=\frac ca=\frac{\sqrt3}{3},\]
\Needspace{5\baselineskip}
故
\nopagebreak[4]
\[c=\frac{\sqrt3}{3}\times\sqrt3=1,\qquad b^2=a^2-c^2=3-1=2.\]
\Needspace{5\baselineskip}
所以椭圆方程为
\nopagebreak[4]
\[\frac{x^2}{3}+\frac{y^2}{2}=1.\]

选 A。

\gkanswer{答案：A。$\triangle AF_1B$ 的周长为 $4a=4\sqrt3$，得 $a=\sqrt3$；又 $e=\dfrac{\sqrt3}{3}$，得 $c=1$，$b^2=2$。}

\end{answer}
\begin{answer}
\gksolhead{7}{}

\textbf{题目}：

曲线 $y=x\mathrm e^{x-1}$ 在点 $(1,1)$ 处切线的斜率等于\mbox{（　　）}

（A）$2\mathrm e$；（B）$\mathrm e$；（C）$2$；（D）$1$。

\Needspace{4\baselineskip}
\textbf{解析}

\Needspace{5\baselineskip}
用乘积求导法则：
\nopagebreak[4]
\[\begin{aligned}
y'&=\left(x\mathrm e^{x-1}\right)'
=\mathrm e^{x-1}+x\mathrm e^{x-1}\\
&=(1+x)\mathrm e^{x-1}.
\end{aligned}\]
\Needspace{5\baselineskip}
在 $x=1$ 处
\nopagebreak[4]
\[y'=(1+1)\mathrm e^{0}=2.\]
故切线斜率为 $2$。

选 C。

\gkanswer{答案：C。$y'=(1+x)\mathrm e^{x-1}$，故 $x=1$ 处斜率为 $2$。}

\end{answer}
\begin{answer}
\gksolhead{8}{}

\textbf{题目}：

正四棱锥的顶点都在同一球面上，若该棱锥的高为 $4$，底面边长为 $2$，则该球的表面积为\mbox{（　　）}

（A）$\dfrac{81\pi}{4}$；（B）$16\pi$；（C）$9\pi$；（D）$\dfrac{27\pi}{4}$。

\Needspace{4\baselineskip}
\textbf{解析}

\Needspace{5\baselineskip}
正四棱锥底面是边长为 $2$ 的正方形，其外接圆半径为
\nopagebreak[4]
\[r=\frac{2\sqrt2}{2}=\sqrt2.\]
\Needspace{5\baselineskip}
球心在棱锥的高所在直线上。设球心到底面的距离为 $t$（$0<t<4$），则球半径
\nopagebreak[4]
\[R=4-t,\]
\Needspace{5\baselineskip}
同时由底面外接圆得
\nopagebreak[4]
\[R^2=t^2+r^2=t^2+2.\]
\Needspace{5\baselineskip}
联立：
\nopagebreak[4]
\[(4-t)^2=t^2+2,\qquad 16-8t+t^2=t^2+2,\]
\Needspace{5\baselineskip}
解得
\nopagebreak[4]
\[t=\frac74,\qquad R=4-\frac74=\frac94.\]
\Needspace{5\baselineskip}
于是球的表面积
\nopagebreak[4]
\[S=4\pi R^2=4\pi\times\frac{81}{16}=\frac{81\pi}{4}.\]

选 A。

\gkanswer{答案：A。设球心到底面距离为 $t$，则 $R^2=t^2+2$ 且 $R=4-t$，解得 $R=\dfrac94$，表面积 $4\pi R^2=\dfrac{81\pi}{4}$。}

\end{answer}
\begin{answer}
\gksolhead{9}{}

\textbf{题目}：

已知双曲线 $C$ 的离心率为 $2$，焦点为 $F_1$、$F_2$，点 $A$ 在 $C$ 上，若 $|F_1A|=2|F_2A|$，则 $\cos\angle AF_2F_1=$\mbox{（　　）}

（A）$\dfrac14$；（B）$\dfrac13$；（C）$\dfrac{\sqrt2}{4}$；（D）$\dfrac{\sqrt2}{3}$。

\Needspace{4\baselineskip}
\textbf{解析}

\Needspace{5\baselineskip}
由离心率 $2$ 得 $c=2a$，故焦距
\nopagebreak[4]
\[|F_1F_2|=2c=4a.\]
\Needspace{5\baselineskip}
设 $A$ 在双曲线的右支上，由双曲线定义
\nopagebreak[4]
\[|AF_1|-|AF_2|=2a.\]
\Needspace{5\baselineskip}
又 $|AF_1|=2|AF_2|$，代入得
\nopagebreak[4]
\[2|AF_2|-|AF_2|=2a,\qquad |AF_2|=2a,\qquad |AF_1|=4a.\]
\Needspace{5\baselineskip}
在 $\triangle AF_1F_2$ 中，对 $\angle AF_2F_1$ 用余弦定理：
\nopagebreak[4]
\[\begin{aligned}
\cos\angle AF_2F_1
&=\frac{|AF_2|^2+|F_1F_2|^2-|AF_1|^2}{2|AF_2|\cdot|F_1F_2|}\\
&=\frac{(2a)^2+(4a)^2-(4a)^2}{2\cdot2a\cdot4a}\\
&=\frac{4a^2}{16a^2}=\frac14.
\end{aligned}\]

选 A。

\gkanswer{答案：A。由定义 $|AF_1|-|AF_2|=2a$ 结合 $|AF_1|=2|AF_2|$ 得 $|AF_2|=2a$、$|AF_1|=4a$，又 $|F_1F_2|=2c=4a$，故 $\cos\angle AF_2F_1=\dfrac{2^2+4^2-4^2}{2\times2\times4}=\dfrac14$。}

\end{answer}
\begin{answer}
\gksolhead{10}{}

\textbf{题目}：

等比数列 $\{a_n\}$ 中，$a_4=2$，$a_5=5$，则数列 $\{\lg a_n\}$ 的前 $8$ 项和等于\mbox{（　　）}

（A）$6$；（B）$5$；（C）$4$；（D）$3$。

\Needspace{4\baselineskip}
\textbf{解析}

\Needspace{5\baselineskip}
设公比为 $q$，则
\nopagebreak[4]
\[q=\frac{a_5}{a_4}=\frac52.\]
于是 $\lg a_n=\lg a_1+(n-1)\lg q$，即 $\{\lg a_n\}$ 是公差为 $\lg q$ 的等差数列。

\Needspace{5\baselineskip}
更简便地利用等比数列的性质：下标和相等的两项之积相等，故
\nopagebreak[4]
\[a_1a_8=a_2a_7=a_3a_6=a_4a_5=2\times5=10.\]
\Needspace{5\baselineskip}
于是
\nopagebreak[4]
\[\begin{aligned}
\sum_{k=1}^{8}\lg a_k
&=\lg(a_1a_2\cdots a_8)\\
&=\lg\left[(a_1a_8)^4\right]
=\lg\left(10^4\right)=4.
\end{aligned}\]

选 C。

\gkanswer{答案：C。$a_1a_8=a_4a_5=10$，故 $\lg a_1+\lg a_8=1$，前 $8$ 项和为 $4$。}

\end{answer}
\begin{answer}
\gksolhead{11}{}

\textbf{题目}：

已知二面角 $\alpha-l-\beta$ 为 $60^\circ$，$AB\subset\alpha$，$AB\perp l$，$A$ 为垂足，$CD\subset\beta$，$C\in l$，$\angle ACD=135^\circ$，则异面直线 $AB$ 与 $CD$ 所成角的余弦值为\mbox{（　　）}

（A）$\dfrac14$；（B）$\dfrac{\sqrt2}{4}$；（C）$\dfrac{\sqrt3}{4}$；（D）$\dfrac12$。

\Needspace{4\baselineskip}
\textbf{解析}

\Needspace{5\baselineskip}
以 $l$ 为 $x$ 轴，$A$ 为原点，平面 $\alpha$ 为 $xOy$ 平面，并取 $AB$ 沿 $y$ 轴正方向：
\nopagebreak[4]
\[A(0,0,0),\qquad \boldsymbol u_{AB}=(0,1,0).\]
\Needspace{5\baselineskip}
二面角 $\alpha-l-\beta$ 为 $60^\circ$，故 $\beta$ 内垂直于 $l$ 的方向可取
\nopagebreak[4]
\[\boldsymbol u=(0,\cos60^\circ,\sin60^\circ)=\left(0,\frac12,\frac{\sqrt3}{2}\right).\]
不妨设 $C(c,0,0)$ 且 $c>0$（$c<0$ 时整体关于 $yOz$ 平面作镜像，夹角不变），则 $\overrightarrow{CA}=(-c,0,0)$，其单位方向为 $(-1,0,0)$。

\Needspace{5\baselineskip}
由 $\angle ACD=135^\circ$，且 $CD\subset\beta$，可把 $\overrightarrow{CD}$ 的方向写成
\nopagebreak[4]
\[\begin{aligned}
\boldsymbol u_{CD}
&=\cos135^\circ\cdot(-1,0,0)+\sin135^\circ\cdot\boldsymbol u\\
&=\frac{\sqrt2}{2}(1,0,0)+\frac{\sqrt2}{2}\left(0,\frac12,\frac{\sqrt3}{2}\right)\\
&=\left(\frac{\sqrt2}{2},\frac{\sqrt2}{4},\frac{\sqrt6}{4}\right).
\end{aligned}\]
\Needspace{5\baselineskip}
于是
\nopagebreak[4]
\[\begin{aligned}
\cos\langle AB,CD\rangle
&=\frac{|\boldsymbol u_{AB}\cdot\boldsymbol u_{CD}|}
{|\boldsymbol u_{AB}|\cdot|\boldsymbol u_{CD}|}\\
&=\frac{\frac{\sqrt2}{4}}{1\cdot1}=\frac{\sqrt2}{4},
\end{aligned}\]
其中 $|\boldsymbol u_{CD}|^2=\dfrac12+\dfrac18+\dfrac38=1$。

选 B。

\gkanswer{答案：B。以 $l$ 为 $x$ 轴、$A$ 为原点、$AB$ 为 $y$ 轴方向建系，得 $\cos\langle AB,CD\rangle=\dfrac{\sqrt2}{4}$。}

\end{answer}
\begin{answer}
\gksolhead{12}{}

\textbf{题目}：

函数 $y=f(x)$ 的图象与函数 $y=g(x)$ 的图象关于直线 $x+y=0$ 对称，则 $y=f(x)$ 的反函数是\mbox{（　　）}

（A）$y=g(x)$；（B）$y=g(-x)$；（C）$y=-g(x)$；（D）$y=-g(-x)$。

\Needspace{4\baselineskip}
\textbf{解析}

\Needspace{5\baselineskip}
点 $(u,v)$ 关于直线 $x+y=0$ 的对称点为 $(-v,-u)$。因此
\nopagebreak[4]
\[v=f(u)\iff -u=g(-v).\]
\Needspace{5\baselineskip}
记 $x=f(u)$，则 $u=f^{-1}(x)$，上式化为
\nopagebreak[4]
\[g(-x)=-f^{-1}(x),\]
\Needspace{5\baselineskip}
即
\nopagebreak[4]
\[f^{-1}(x)=-g(-x).\]
所以 $y=f(x)$ 的反函数是 $y=-g(-x)$。

选 D。

\gkanswer{答案：D。关于 $x+y=0$ 对称给出 $g(x)=-f^{-1}(-x)$，故 $f^{-1}(x)=-g(-x)$。}

\end{answer}
\begin{answer}
\gksolhead{13}{}

\textbf{题目}：

$\left(\dfrac{x}{\sqrt y}-\dfrac{y}{\sqrt x}\right)^8$ 的展开式中 $x^2y^2$ 的系数为\gkblank{}。（用数字作答）

\Needspace{4\baselineskip}
\textbf{解析}

\Needspace{5\baselineskip}
把两项分别写成幂：
\nopagebreak[4]
\[\frac{x}{\sqrt y}=xy^{-\frac12},\qquad \frac{y}{\sqrt x}=yx^{-\frac12}.\]
\Needspace{5\baselineskip}
二项展开的通项为
\nopagebreak[4]
\[\begin{aligned}
T_{k+1}
&=\mathrm C_8^k\left(xy^{-\frac12}\right)^{8-k}\left(-yx^{-\frac12}\right)^k\\
&=\mathrm C_8^k(-1)^k x^{8-k-\frac k2}y^{-\frac{8-k}{2}+k}\\
&=\mathrm C_8^k(-1)^k x^{8-\frac{3k}{2}}y^{\frac{3k}{2}-4}.
\end{aligned}\]
\Needspace{5\baselineskip}
要求 $x^2y^2$ 项，令
\nopagebreak[4]
\[8-\frac{3k}{2}=2\ \Longrightarrow\ k=4,\]
\Needspace{5\baselineskip}
此时 $y$ 的指数为 $\dfrac{3\times4}{2}-4=2$，一致。故该项系数为
\nopagebreak[4]
\[\mathrm C_8^4(-1)^4=70.\]

\gkanswer{答案：$70$。通项为 $\mathrm C_8^k(-1)^k x^{8-\frac{3k}{2}}y^{\frac{3k}{2}-4}$，令指数为 $2$ 得 $k=4$，系数 $\mathrm C_8^4=70$。}

\end{answer}
\begin{answer}
\gksolhead{14}{}

\textbf{题目}：

\Needspace{5\baselineskip}
设 $x$、$y$ 满足约束条件
\nopagebreak[4]
\[\begin{cases}
x-y\geqslant 0,\\
x+2y\leqslant 3,\\
x-2y\leqslant 1,
\end{cases}\]
则 $z=x+4y$ 的最大值为\gkblank{}。

\Needspace{4\baselineskip}
\textbf{解析}

\Needspace{5\baselineskip}
画出可行域。三条边界两两相交，得三个顶点：
\nopagebreak[4]
\[\begin{cases}x=y,\\ x+2y=3\end{cases}\Longrightarrow (1,1);\qquad
\begin{cases}x=y,\\ x-2y=1\end{cases}\Longrightarrow (-1,-1);\]
\nopagebreak[4]
\[\begin{cases}x+2y=3,\\ x-2y=1\end{cases}\Longrightarrow \left(2,\frac12\right).\]
\Needspace{5\baselineskip}
由于 $z=x+4y$ 是线性函数，最大值必在顶点处取得。逐点计算：
\nopagebreak[4]
\[z(1,1)=1+4=5,\qquad z(-1,-1)=-1-4=-5,\qquad z\left(2,\tfrac12\right)=2+2=4.\]
故最大值为 $5$。

\gkanswer{答案：$5$。可行域的三个顶点为 $(1,1)$、$(-1,-1)$、$\left(2,\dfrac12\right)$，对应 $z$ 值为 $5$、$-5$、$4$，最大值为 $5$。}

\end{answer}
\begin{answer}
\gksolhead{15}{}

\textbf{题目}：

直线 $l_1$ 和 $l_2$ 是圆 $x^2+y^2=2$ 的两条切线，若 $l_1$ 与 $l_2$ 的交点为 $(1,3)$，则 $l_1$ 与 $l_2$ 的夹角的正切值等于\gkblank{}。

\Needspace{4\baselineskip}
\textbf{解析}

\Needspace{5\baselineskip}
设两切线的夹角为 $\theta$，切点为 $T$。交点 $P(1,3)$ 到圆心 $O(0,0)$ 的距离为
\nopagebreak[4]
\[|PO|=\sqrt{1^2+3^2}=\sqrt{10},\]
\Needspace{5\baselineskip}
圆半径 $r=\sqrt2$。由切线性质 $OT\perp PT$，在直角三角形 $POT$ 中
\nopagebreak[4]
\[\sin\frac{\theta}{2}=\frac{r}{|PO|}=\frac{\sqrt2}{\sqrt{10}}=\frac1{\sqrt5}.\]
\Needspace{5\baselineskip}
于是
\nopagebreak[4]
\[\cos\frac{\theta}{2}=\sqrt{1-\frac15}=\frac{2}{\sqrt5},\qquad
\tan\frac{\theta}{2}=\frac{1}{2}.\]
\Needspace{5\baselineskip}
由二倍角公式
\nopagebreak[4]
\[\tan\theta=\frac{2\tan\frac{\theta}{2}}{1-\tan^2\frac{\theta}{2}}
=\frac{2\times\frac12}{1-\frac14}=\frac43.\]

\gkanswer{答案：$\dfrac43$。设夹角为 $\theta$，则 $\sin\dfrac{\theta}{2}=\dfrac{\sqrt2}{\sqrt{10}}=\dfrac1{\sqrt5}$，$\tan\dfrac{\theta}{2}=\dfrac12$，故 $\tan\theta=\dfrac43$。}

\end{answer}
\begin{answer}
\gksolhead{16}{}

\textbf{题目}：

若函数 $f(x)=\cos2x+a\sin x$ 在区间 $\left(\dfrac\pi6,\dfrac\pi2\right)$ 上是减函数，则 $a$ 的取值范围是\gkblank{}。

\Needspace{4\baselineskip}
\textbf{解析}

\Needspace{5\baselineskip}
求导：
\nopagebreak[4]
\[\begin{aligned}
f'(x)&=-2\sin2x+a\cos x\\
&=\cos x\,(a-4\sin x).
\end{aligned}\]
\Needspace{5\baselineskip}
在 $\left(\dfrac\pi6,\dfrac\pi2\right)$ 上 $\cos x>0$，故
\nopagebreak[4]
\[f'(x)\leqslant0\iff a-4\sin x\leqslant0\iff a\leqslant4\sin x.\]
\Needspace{5\baselineskip}
要让上式对一切 $x\in\left(\dfrac\pi6,\dfrac\pi2\right)$ 成立，只需
\nopagebreak[4]
\[a\leqslant\inf_{x\in(\frac\pi6,\frac\pi2)}4\sin x=4\sin\frac\pi6=2.\]
当 $a=2$ 时，$a-4\sin x=2(1-2\sin x)<0$ 在区间内严格成立，$f$ 确实递减。

反过来，若 $a>2$，取 $x_0\in\left(\dfrac\pi6,\dfrac\pi2\right)$ 充分接近 $\dfrac\pi6$，使 $4\sin x_0<a$，则 $f'(x_0)>0$，$f$ 在 $x_0$ 附近递增，与“在 $\left(\dfrac\pi6,\dfrac\pi2\right)$ 上是减函数”矛盾。故 $a\leqslant2$ 也是必要的。

所以 $a$ 的取值范围是 $(-\infty,2]$。

\gkanswer{答案：$(-\infty,2]$。在 $\left(\dfrac\pi6,\dfrac\pi2\right)$ 上 $\cos x>0$，$f'(x)\leqslant0$ 等价于 $a\leqslant4\sin x$ 恒成立，右端下确界为 $4\sin\dfrac\pi6=2$。}

\end{answer}
\begin{answer}
\gksolhead{17}{}

\textbf{题目}：

$\triangle ABC$ 的内角 $A$、$B$、$C$ 的对边分别为 $a$、$b$、$c$，已知 $3a\cos C=2c\cos A$，$\tan A=\dfrac13$，求 $B$。

\Needspace{4\baselineskip}
\textbf{解析}

\Needspace{5\baselineskip}
由正弦定理 $\dfrac{a}{\sin A}=\dfrac{c}{\sin C}$，把 $a$、$c$ 换成角：
\nopagebreak[4]
\[3\sin A\cos C=2\sin C\cos A.\]
\Needspace{5\baselineskip}
由 $\tan A=\dfrac13>0$ 知 $A$ 为锐角，$\cos A>0$。两边同除以 $\cos A$：
\nopagebreak[4]
\[3\tan A\cos C=2\sin C.\]
\Needspace{5\baselineskip}
若 $\cos C=0$，则 $\sin C=1$，上式左边为 $0$、右边为 $2$，矛盾，故 $\cos C\neq0$。于是
\nopagebreak[4]
\[3\tan A=2\tan C.\]
\Needspace{5\baselineskip}
由 $\tan A=\dfrac13$ 得
\nopagebreak[4]
\[\tan C=\frac32\times\frac13=\frac12,\]
从而 $C$ 也是锐角。
\Needspace{5\baselineskip}
于是
\nopagebreak[4]
\[\tan(A+C)=\frac{\tan A+\tan C}{1-\tan A\tan C}
=\frac{\frac13+\frac12}{1-\frac16}=\frac{\frac56}{\frac56}=1.\]
\Needspace{5\baselineskip}
又 $A+C\in(0,\pi)$，且 $\tan(A+C)>0$，故
\nopagebreak[4]
\[A+C=\frac\pi4,\]
\Needspace{5\baselineskip}
从而
\nopagebreak[4]
\[B=\pi-(A+C)=\frac{3\pi}{4}.\]

\gkanswer{答案：$B=\dfrac{3\pi}{4}$。由正弦定理得 $3\sin A\cos C=2\sin C\cos A$，即 $3\tan A=2\tan C$，故 $\tan C=\dfrac12$，$\tan(A+C)=1$，$A+C=\dfrac{\pi}{4}$。}

\end{answer}
\begin{answer}
\gksolhead{18}{}

\textbf{题目}：

等差数列 $\{a_n\}$ 的前 $n$ 项和为 $S_n$，已知 $a_1=10$，$a_2$ 为整数，且 $S_n\leqslant S_4$。

（1） 求 $\{a_n\}$ 的通项公式；

（2） 设 $b_n=\dfrac{1}{a_na_{n+1}}$，求数列 $\{b_n\}$ 的前 $n$ 项和 $T_n$。

\Needspace{4\baselineskip}
\textbf{解析}

\Needspace{5\baselineskip}
设公差为 $d$，则 $a_n=10+(n-1)d$，且
\nopagebreak[4]
\[S_n=10n+\frac{d\,n(n-1)}{2}.\]

\gksubpart{（1）}

\Needspace{5\baselineskip}
当 $d<0$ 时 $S_n$ 的图象是开口向下的抛物线，$S_n\leqslant S_4$ 对一切正整数 $n$ 成立，等价于
\nopagebreak[4]
\[a_4\geqslant0,\qquad a_5\leqslant0.\]
\Needspace{5\baselineskip}
即
\nopagebreak[4]
\[10+3d\geqslant0,\qquad 10+4d\leqslant0,\]
\Needspace{5\baselineskip}
所以
\nopagebreak[4]
\[-\frac{10}{3}\leqslant d\leqslant-\frac52.\]
\Needspace{5\baselineskip}
又 $a_2=10+d$ 为整数，故 $d=a_2-10$ 是整数。落在上述区间内的整数只有
\nopagebreak[4]
\[d=-3.\]
\Needspace{5\baselineskip}
于是
\nopagebreak[4]
\[a_n=10-3(n-1)=13-3n.\]
此时 $a_4=1>0$，$a_5=-2<0$，$S_n$ 确在 $n=4$ 处取最大值。

\gksubpart{（2）}

\Needspace{5\baselineskip}
由（1）得
\nopagebreak[4]
\[b_n=\frac{1}{a_na_{n+1}}=\frac{1}{(13-3n)(10-3n)}.\]
\Needspace{5\baselineskip}
注意分母中两因子之差为
\nopagebreak[4]
\[(13-3n)-(10-3n)=3,\]
\Needspace{5\baselineskip}
故可裂项：
\nopagebreak[4]
\[b_n=\frac13\left(\frac{1}{10-3n}-\frac{1}{13-3n}\right).\]
\Needspace{5\baselineskip}
求和：
\nopagebreak[4]
\[\begin{aligned}
T_n&=\frac13\sum_{k=1}^{n}\left(\frac{1}{10-3k}-\frac{1}{13-3k}\right).
\end{aligned}\]
\Needspace{5\baselineskip}
第二个和对下标平移一位，即 $\dfrac{1}{13-3k}=\dfrac{1}{10-3(k-1)}$，于是两和相消，只剩下
\nopagebreak[4]
\[\begin{aligned}
T_n&=\frac13\left(\frac{1}{10-3n}-\frac{1}{10}\right)\\
&=\frac13\cdot\frac{10-(10-3n)}{10(10-3n)}
=\frac{n}{10(10-3n)}.
\end{aligned}\]
例如 $n=1$ 时 $T_1=\dfrac{1}{70}=b_1$，与直接计算一致。

\gkanswer{答案（1）：$a_n=13-3n$。由 $S_n\leqslant S_4$ 得 $a_4\geqslant0$、$a_5\leqslant0$，即 $-\dfrac{10}{3}\leqslant d\leqslant-\dfrac52$，又 $a_2=10+d$ 为整数，故 $d=-3$。}

\gkanswer{答案（2）：$T_n=\dfrac{n}{10(10-3n)}$。由 $b_n=\dfrac13\left(\dfrac1{10-3n}-\dfrac1{13-3n}\right)$ 裂项相消得 $T_n=\dfrac13\left(\dfrac1{10-3n}-\dfrac1{10}\right)$。}

\end{answer}
\begin{answer}
\gksolhead{19}{}

\textbf{题目}：

如图，三棱柱 $ABC-A_1B_1C_1$ 中，点 $A_1$ 在平面 $ABC$ 内的射影 $D$ 在 $AC$ 上，$\angle ACB=90^\circ$，$BC=1$，$AC=CC_1=2$。
\nopagebreak[4]
\begin{center}% ============================================================
% f2014_outline_science_q19_prism.tex
% 2014 全国大纲卷（理科）第 19 题：三棱柱 ABC-A1B1C1
% A1 在平面 ABC 内的射影 D 在 AC 上，∠ACB=90°。
% 来源 PDF：普通高考/2014/2014大纲理(广西).pdf
% 原页：第 1 页（A3 横排）第 3 栏，印面 869，第 19 题题干下方
% 画法为平行投影示意，不按比例；虚线为不可见棱及辅助线。
% ============================================================
\begin{tikzpicture}[>=Stealth, line width=0.7pt, font=\small,
  line cap=round, line join=round]
  % 底面 ABC：C 为直角顶点，CA=2、CB=1；D 为 AC 中点，
  % A1 在 D 正上方（h=√3），三棱柱侧棱 AA1∥CC1∥BB1。
  \coordinate (A)  at (0.00,0.00);
  \coordinate (C)  at (1.80,1.15);
  \coordinate (B)  at (3.90,1.00);
  \coordinate (D)  at (0.90,0.58);
  \coordinate (A1) at (0.95,3.38);
  \coordinate (C1) at (2.75,4.53);
  \coordinate (B1) at (4.85,4.38);

  % 可见棱
  \draw (A) -- (B) -- (B1) -- (C1) -- (A1) -- cycle;
  \draw (A) -- (A1);
  \draw (A1) -- (B1);
  \draw (A1) -- (B);

  % 不可见棱与辅助线
  \draw[dashed] (A) -- (C) -- (B);
  \draw[dashed] (C) -- (C1);
  \draw[dashed] (A1) -- (D);
  \draw[dashed] (A1) -- (C);

  \fill (A) circle (1.1pt) node[below left] {$A$};
  \fill (B) circle (1.1pt) node[right] {$B$};
  \fill (C) circle (1.1pt) node[above left] {$C$};
  \fill (D) circle (1.1pt) node[below=1pt] {$D$};
  \fill (A1) circle (1.1pt) node[left] {$A_1$};
  \fill (C1) circle (1.1pt) node[above] {$C_1$};
  \fill (B1) circle (1.1pt) node[above right] {$B_1$};
\end{tikzpicture}
\end{center}
\nopagebreak[4]
（1） 证明：$AC_1\perp A_1B$；

（2） 设直线 $AA_1$ 与平面 $BCC_1B_1$ 的距离为 $\sqrt3$，求二面角 $A_1-AB-C$ 的大小。

\Needspace{4\baselineskip}
\textbf{解析}

\Needspace{5\baselineskip}
以 $C$ 为原点，$CA$ 所在直线为 $x$ 轴，$CB$ 所在直线为 $y$ 轴建立空间直角坐标系，则
\nopagebreak[4]
\[C(0,0,0),\quad A(2,0,0),\quad B(0,1,0).\]
\Needspace{5\baselineskip}
设 $D(t,0,0)$，则 $A_1(t,0,h)$，其中 $0\leqslant t\leqslant2$，$h>0$。因为 $CC_1\parallel AA_1$ 且 $CC_1=AA_1=2$，所以
\nopagebreak[4]
\[(t-2)^2+h^2=4,\]
\Needspace{5\baselineskip}
即
\nopagebreak[4]
\[h^2=4t-t^2. \tag{1}\]

\gksubpart{（1）}

\Needspace{5\baselineskip}
由 $\overrightarrow{CC_1}=\overrightarrow{AA_1}=(t-2,0,h)$，得
\nopagebreak[4]
\[C_1(t-2,0,h).\]
\Needspace{5\baselineskip}
于是
\nopagebreak[4]
\[\overrightarrow{AC_1}=(t-4,0,h),\qquad
\overrightarrow{A_1B}=(-t,1,-h).\]
\Needspace{5\baselineskip}
计算内积并利用式（1）：
\nopagebreak[4]
\[\begin{aligned}
\overrightarrow{AC_1}\cdot\overrightarrow{A_1B}
&=(t-4)(-t)+0\cdot1+h(-h)\\
&=-t^2+4t-h^2\\
&=4t-t^2-h^2=0.
\end{aligned}\]
\Needspace{5\baselineskip}
故
\nopagebreak[4]
\[AC_1\perp A_1B.\]

\gksubpart{（2）}

\Needspace{5\baselineskip}
平面 $BCC_1B_1$ 由 $C$、$B$、$C_1$ 确定。取
\nopagebreak[4]
\[\boldsymbol n=\overrightarrow{CB}\times\overrightarrow{CC_1}
=(0,1,0)\times(t-2,0,h)=(h,0,2-t),\]
\Needspace{5\baselineskip}
则 $\boldsymbol n$ 是该平面的法向量，且
\nopagebreak[4]
\[|\boldsymbol n|=\sqrt{h^2+(2-t)^2}.\]
\Needspace{5\baselineskip}
点 $A$ 到该平面的距离为
\nopagebreak[4]
\[d=\frac{|\boldsymbol n\cdot\overrightarrow{CA}|}{|\boldsymbol n|}
=\frac{|2h|}{\sqrt{h^2+(2-t)^2}}.\]
\Needspace{5\baselineskip}
已知 $d=\sqrt3$，故
\nopagebreak[4]
\[\frac{4h^2}{h^2+(2-t)^2}=3,\]
\Needspace{5\baselineskip}
即
\nopagebreak[4]
\[h^2=3(2-t)^2. \tag{2}\]
\Needspace{5\baselineskip}
把式（1）代入式（2）：
\nopagebreak[4]
\[4t-t^2=3(4-4t+t^2),\]
\Needspace{5\baselineskip}
整理得
\nopagebreak[4]
\[4t^2-16t+12=0,\qquad t^2-4t+3=0,\]
\Needspace{5\baselineskip}
所以
\nopagebreak[4]
\[t=1\quad\text{\text{或}}\quad t=3.\]
\Needspace{5\baselineskip}
因为 $D$ 在线段 $AC$ 上，$0\leqslant t\leqslant2$，故
\nopagebreak[4]
\[t=1,\qquad h=\sqrt3.\]
\Needspace{5\baselineskip}
于是
\nopagebreak[4]
\[A_1(1,0,\sqrt3).\]

\Needspace{5\baselineskip}
平面 $A_1AB$ 的法向量为
\nopagebreak[4]
\[\begin{aligned}
\boldsymbol m&=\overrightarrow{AB}\times\overrightarrow{AA_1}\\
&=(-2,1,0)\times(-1,0,\sqrt3)\\
&=(\sqrt3,\,2\sqrt3,\,1).
\end{aligned}\]
\Needspace{5\baselineskip}
平面 $ABC$ 的法向量可取 $\boldsymbol n_0=(0,0,1)$。设二面角 $A_1-AB-C$ 的大小为 $\theta$，则
\nopagebreak[4]
\[\cos\theta
=\frac{|\boldsymbol m\cdot\boldsymbol n_0|}{|\boldsymbol m|\,|\boldsymbol n_0|}
=\frac{1}{\sqrt{3+12+1}\cdot1}
=\frac14.\]
\Needspace{5\baselineskip}
结合图形知该二面角为锐角，所以
\nopagebreak[4]
\[\theta=\arccos\frac14.\]

\gkanswer{答案（1）：见解析版}

\gkanswer{答案（2）：$\arccos\dfrac14$。由线面距离 $\sqrt3$ 得 $h^2=3(2-t)^2$，结合 $h^2=4t-t^2$ 及 $D$ 在 $AC$ 上得 $t=1$、$h=\sqrt3$，两半平面法向量夹角的余弦为 $\dfrac14$。}

\end{answer}
\begin{answer}
\gksolhead{20}{}

\textbf{题目}：

设每个工作日甲、乙、丙、丁 4 人需使用某种设备的概率分别为 $0.6$、$0.5$、$0.5$、$0.4$，各人是否需使用设备相互独立。

（1） 求同一工作日至少 3 人需使用设备的概率；

（2） $X$ 表示同一工作日需使用设备的人数，求 $X$ 的数学期望。

\Needspace{4\baselineskip}
\textbf{解析}

设同一工作日需使用设备的人数为 $X$，四人是否使用设备相互独立。

\gksubpart{（1）}

“至少 3 人使用”即 $X=3$ 或 $X=4$。

\Needspace{5\baselineskip}
$X=4$ 的概率为
\nopagebreak[4]
\[P(X=4)=0.6\times0.5\times0.5\times0.4=0.06.\]
\Needspace{5\baselineskip}
$X=3$ 即恰有一人不用，按“不用者是谁”分类：
\nopagebreak[4]
\[\begin{aligned}
P(X=3)&=0.4\times0.5\times0.5\times0.4
+0.6\times0.5\times0.5\times0.4\\
&\quad+0.6\times0.5\times0.5\times0.4
+0.6\times0.5\times0.5\times0.6\\
&=0.04+0.06+0.06+0.09=0.25.
\end{aligned}\]
\Needspace{5\baselineskip}
所以
\nopagebreak[4]
\[P(X\geqslant3)=P(X=3)+P(X=4)=0.25+0.06=0.31.\]

\gksubpart{（2）}

\Needspace{5\baselineskip}
引入示性变量 $X_i$：第 $i$ 人使用设备时 $X_i=1$，否则 $X_i=0$。四人是否使用设备相互独立，且 $E(X_i)$ 就是该人使用设备的概率，故
\nopagebreak[4]
\[E(X_i)=p_i,\qquad i=1,2,3,4,\]
\Needspace{5\baselineskip}
其中 $p_1=0.6$，$p_2=p_3=0.5$，$p_4=0.4$。由 $X=X_1+X_2+X_3+X_4$ 及期望的线性性，
\nopagebreak[4]
\[\begin{aligned}
E(X)&=E(X_1)+E(X_2)+E(X_3)+E(X_4)\\
&=0.6+0.5+0.5+0.4=2.
\end{aligned}\]
即 $X$ 的数学期望为 $2$。

（附：由分布列 $P(X=0)=0.06$、$P(X=1)=0.25$、$P(X=2)=0.38$、$P(X=3)=0.25$、$P(X=4)=0.06$ 求和也得 $E(X)=2$。）

\gkanswer{答案（1）：$0.31$。$P(X=4)=0.06$，$P(X=3)=0.25$，故 $P(X\geqslant3)=0.31$。}

\gkanswer{答案（2）：$2$。四人是否使用设备相互独立，期望等于各人概率之和 $0.6+0.5+0.5+0.4=2$。}

\end{answer}
\begin{answer}
\gksolhead{21}{}

\textbf{题目}：

已知抛物线 $C:y^2=2px$（$p>0$）的焦点为 $F$，直线 $y=4$ 与 $y$ 轴的交点为 $P$，与 $C$ 的交点为 $Q$，且 $|QF|=\dfrac54|PQ|$。

（1） 求 $C$ 的方程；

（2） 过 $F$ 的直线 $l$ 与 $C$ 相交于 $A$、$B$ 两点，若 $AB$ 的垂直平分线 $l'$ 与 $C$ 相交于 $M$、$N$ 两点，且 $A$、$M$、$B$、$N$ 四点在同一圆上，求 $l$ 的方程。

\Needspace{4\baselineskip}
\textbf{解析}

\gksubpart{（1）}

\Needspace{5\baselineskip}
抛物线 $C:y^2=2px$ 的焦点为
\nopagebreak[4]
\[F\left(\frac p2,0\right).\]
\Needspace{5\baselineskip}
直线 $y=4$ 与 $y$ 轴交于
\nopagebreak[4]
\[P(0,4).\]
\Needspace{5\baselineskip}
与 $C$ 交于 $Q$，由 $16=2px$ 得
\nopagebreak[4]
\[x_Q=\frac8p,\qquad Q\left(\frac8p,4\right).\]
\Needspace{5\baselineskip}
于是
\nopagebreak[4]
\[|PQ|=\frac8p.\]
\Needspace{5\baselineskip}
对抛物线上的点 $Q$，到焦点 $F$ 的距离等于到准线 $x=-\dfrac p2$ 的距离，故
\nopagebreak[4]
\[|QF|=x_Q+\frac p2=\frac8p+\frac p2.\]
\Needspace{5\baselineskip}
由 $|QF|=\dfrac54|PQ|$ 得
\nopagebreak[4]
\[\frac8p+\frac p2=\frac54\cdot\frac8p=\frac{10}{p}.\]
\Needspace{5\baselineskip}
两边乘以 $p$：
\nopagebreak[4]
\[8+\frac{p^2}{2}=10,\qquad p^2=4.\]
\Needspace{5\baselineskip}
因为 $p>0$，所以 $p=2$。故
\nopagebreak[4]
\[C:y^2=4x.\]

\gksubpart{（2）}

\Needspace{5\baselineskip}
此时 $F(1,0)$。设直线 $l$ 的方程为
\nopagebreak[4]
\[x=my+1.\]
若 $m=0$，则 $l:x=1$，$AB$ 的垂直平分线为 $y=0$，它与 $C$ 只有公共点 $(0,0)$，不能交出 $M$、$N$ 两点，故 $m\neq0$。

\Needspace{5\baselineskip}
设 $A(x_1,y_1)$、$B(x_2,y_2)$。把 $x=my+1$ 代入 $y^2=4x$：
\nopagebreak[4]
\[y^2-4my-4=0.\]
\Needspace{5\baselineskip}
由韦达定理
\nopagebreak[4]
\[y_1+y_2=4m,\qquad y_1y_2=-4.\]
\Needspace{5\baselineskip}
$AB$ 的中点坐标为
\nopagebreak[4]
\[\left(\frac{x_1+x_2}{2},\frac{y_1+y_2}{2}\right)
=(2m^2+1,\,2m).\]
\Needspace{5\baselineskip}
直线 $l$ 的斜率为 $\dfrac1m$，故 $l'$ 的斜率为 $-m$，其方程为
\nopagebreak[4]
\[y-2m=-m(x-2m^2-1),\]
\Needspace{5\baselineskip}
即
\nopagebreak[4]
\[y=-mx+2m^3+3m.\]
\Needspace{5\baselineskip}
设 $M(x_3,y_3)$、$N(x_4,y_4)$。由 $l'$ 的方程解出
\nopagebreak[4]
\[x=2m^2+3-\frac ym,\]
\Needspace{5\baselineskip}
代入 $y^2=4x$：
\nopagebreak[4]
\[y^2=8m^2+12-\frac{4y}{m},\]
\Needspace{5\baselineskip}
即
\nopagebreak[4]
\[y^2+\frac4m y-(8m^2+12)=0.\]
\Needspace{5\baselineskip}
故
\nopagebreak[4]
\[y_3+y_4=-\frac4m.\]

\Needspace{5\baselineskip}
下面利用四点共圆的条件。任取一个圆
\nopagebreak[4]
\[x^2+y^2+Dx+Ey+F=0.\]
\Needspace{5\baselineskip}
把 $x=\dfrac{y^2}{4}$ 代入，得
\nopagebreak[4]
\[\frac1{16}y^4+\left(1+\frac D4\right)y^2+Ey+F=0.\]
这个四次方程的四个根正是圆与抛物线交点的纵坐标。其 $y^3$ 项系数为 $0$，故四根之和为 $0$。反过来，若抛物线上四点的纵坐标之和为 $0$，则以它们为根的首一四次多项式没有 $y^3$ 项，适当选取 $D$、$E$、$F$ 即可写成上式，从而四点共圆。

\Needspace{5\baselineskip}
因此 $A$、$M$、$B$、$N$ 四点共圆等价于
\nopagebreak[4]
\[y_1+y_2+y_3+y_4=0.\]
\Needspace{5\baselineskip}
代入上面的和：
\nopagebreak[4]
\[4m-\frac4m=0,\]
\Needspace{5\baselineskip}
故
\nopagebreak[4]
\[m^2=1,\qquad m=\pm1.\]
当 $m=1$ 时，$l:x=y+1$，即 $y=x-1$；当 $m=-1$ 时，$l:x=-y+1$，即 $y=-x+1$。此时 $l'$ 与 $C$ 均有两个不同交点，四点共圆的条件确实成立。

\Needspace{5\baselineskip}
所以 $l$ 的方程为
\nopagebreak[4]
\[y=x-1\quad\text{\text{或}}\quad y=-x+1.\]

\gkanswer{答案（1）：$y^2=4x$。由 $|PQ|=\dfrac8p$、$|QF|=\dfrac8p+\dfrac p2$ 及 $|QF|=\dfrac54|PQ|$ 解得 $p=2$。}

\gkanswer{答案（2）：$l:y=x-1$ 或 $l:y=-x+1$。设 $l:x=my+1$，由 $A$、$M$、$B$、$N$ 四点共圆得四纵坐标之和为 $0$，即 $4m-\dfrac4m=0$，故 $m=\pm1$。}

\end{answer}
\begin{answer}
\gksolhead{22}{}

\textbf{题目}：

函数 $f(x)=\ln(x+1)-\dfrac{ax}{x+a}$（$a>1$）。

（1） 讨论 $f(x)$ 的单调性；

（2） 设 $a_1=1$，$a_{n+1}=\ln(a_n+1)$，证明：$\dfrac{2}{n+2}<a_n\leqslant\dfrac{3}{n+2}$。

\Needspace{4\baselineskip}
\textbf{解析}

\gksubpart{（1）}

\Needspace{5\baselineskip}
定义域为 $(-1,+\infty)$，在此区间上 $x+1>0$，$x+a>a-1>0$。求导：
\nopagebreak[4]
\[\begin{aligned}
f'(x)&=\frac{1}{x+1}-\frac{a(x+a)-ax}{(x+a)^2}\\
&=\frac{1}{x+1}-\frac{a^2}{(x+a)^2}\\
&=\frac{(x+a)^2-a^2(x+1)}{(x+1)(x+a)^2}\\
&=\frac{x^2+(2a-a^2)x}{(x+1)(x+a)^2}\\
&=\frac{x\bigl[x-a(a-2)\bigr]}{(x+1)(x+a)^2}.
\end{aligned}\]
分母恒正，故 $f'(x)$ 的符号由分子 $x\bigl[x-a(a-2)\bigr]$ 决定，两个零点为 $0$ 与 $a(a-2)$。

\Needspace{5\baselineskip}
当 $1<a<2$ 时，$a(a-2)=(a-1)^2-1\in(-1,0)$。于是
\nopagebreak[4]
\[f'(x)>0\ (x\in(-1,a(a-2))),\quad
f'(x)<0\ (x\in(a(a-2),0)),\quad
f'(x)>0\ (x\in(0,+\infty)).\]
所以 $f$ 在 $(-1,a(a-2))$、$(0,+\infty)$ 上递增，在 $(a(a-2),0)$ 上递减。

\Needspace{5\baselineskip}
当 $a=2$ 时，
\nopagebreak[4]
\[f'(x)=\frac{x^2}{(x+1)(x+2)^2}\geqslant0,\]
仅在 $x=0$ 处取零，故 $f$ 在 $(-1,+\infty)$ 上递增。

\Needspace{5\baselineskip}
当 $a>2$ 时，$a(a-2)>0$。于是
\nopagebreak[4]
\[f'(x)>0\ (x\in(-1,0)),\quad
f'(x)<0\ (x\in(0,a(a-2))),\quad
f'(x)>0\ (x\in(a(a-2),+\infty)).\]
所以 $f$ 在 $(-1,0)$、$(a(a-2),+\infty)$ 上递增，在 $(0,a(a-2))$ 上递减。

\gksubpart{（2）}

先用归纳法证明两个基本事实。

其一，$a_n>0$。$a_1=1>0$；若 $a_n>0$，则 $a_{n+1}=\ln(1+a_n)>\ln1=0$。

其二，$a_n\leqslant1$。对 $x>0$ 有 $\ln(1+x)<x$，故 $a_{n+1}=\ln(1+a_n)<a_n$，数列递减，$a_n\leqslant a_1=1$。

\Needspace{5\baselineskip}
由（1），当 $a=2$ 时
\nopagebreak[4]
\[f(x)=\ln(x+1)-\frac{2x}{x+2}\]
\Needspace{5\baselineskip}
在 $(-1,+\infty)$ 上递增，且 $f(0)=0$，故 $x>0$ 时
\nopagebreak[4]
\[\ln(1+x)>\frac{2x}{x+2}.\]
\Needspace{5\baselineskip}
取 $x=a_n>0$，得
\nopagebreak[4]
\[a_{n+1}>\frac{2a_n}{a_n+2}.\]
\Needspace{5\baselineskip}
两边取倒数（各量均为正）：
\nopagebreak[4]
\[\frac{1}{a_{n+1}}<\frac{a_n+2}{2a_n}=\frac{1}{a_n}+\frac12.\]
\Needspace{5\baselineskip}
对 $k=1,2,\cdots,n-1$ 累加，得
\nopagebreak[4]
\[\frac{1}{a_n}<\frac{1}{a_1}+\frac{n-1}{2}=\frac{n+1}{2},\]
\Needspace{5\baselineskip}
故
\nopagebreak[4]
\[a_n>\frac{2}{n+1}>\frac{2}{n+2}\qquad(n\geqslant2),\]
而 $n=1$ 时 $a_1=1>\dfrac23$。左半不等式得证。

\Needspace{5\baselineskip}
再由（1），当 $a=3>2$ 时
\nopagebreak[4]
\[f(x)=\ln(x+1)-\frac{3x}{x+3}\]
\Needspace{5\baselineskip}
在 $(0,3)$ 上递减，且 $f(0)=0$，故 $0<x<3$ 时
\nopagebreak[4]
\[\ln(1+x)<\frac{3x}{x+3}.\]
\Needspace{5\baselineskip}
因 $0<a_n\leqslant1<3$，取 $x=a_n$ 得
\nopagebreak[4]
\[a_{n+1}<\frac{3a_n}{a_n+3},\]
\Needspace{5\baselineskip}
取倒数：
\nopagebreak[4]
\[\frac{1}{a_{n+1}}>\frac{1}{a_n}+\frac13.\]
\Needspace{5\baselineskip}
累加得
\nopagebreak[4]
\[\frac{1}{a_n}>\frac{1}{a_1}+\frac{n-1}{3}=\frac{n+2}{3},\]
\Needspace{5\baselineskip}
故
\nopagebreak[4]
\[a_n<\frac{3}{n+2}\qquad(n\geqslant2),\]
\Needspace{5\baselineskip}
而 $n=1$ 时 $a_1=1=\dfrac33$。所以对一切正整数 $n$ 有
\nopagebreak[4]
\[a_n\leqslant\frac{3}{n+2}.\]
\Needspace{5\baselineskip}
综上，
\nopagebreak[4]
\[\frac{2}{n+2}<a_n\leqslant\frac{3}{n+2}.\]

\gkanswer{答案（1）：当 $1<a<2$ 时在 $(-1,a(a-2))$、$(0,+\infty)$ 上递增，在 $(a(a-2),0)$ 上递减；当 $a=2$ 时在 $(-1,+\infty)$ 上递增；当 $a>2$ 时在 $(-1,0)$、$(a(a-2),+\infty)$ 上递增，在 $(0,a(a-2))$ 上递减。}

\gkanswer{答案（2）：见解析版}

\end{answer}